% concatenated from main.tex
\documentclass[11pt]{amsart}

\usepackage[T1]{fontenc}
\usepackage{lmodern}
\usepackage{microtype}
\usepackage{mathtools}
\usepackage{amssymb}
\usepackage{enumitem}
\usepackage[hidelinks]{hyperref}
\usepackage[nameinlink,noabbrev]{cleveref}

\allowdisplaybreaks
\setlist[enumerate]{label=(\roman*),leftmargin=2.1em}

\newcommand{\N}{\mathbb N}
\newcommand{\Z}{\mathbb Z}
\newcommand{\Q}{\mathbb Q}
\newcommand{\cP}{\mathcal P}
\newcommand{\cI}{\mathcal I}
\newcommand{\fd}{\operatorname{fd}}
\newcommand{\radQ}{\operatorname{rad}_{\Q}}
\newcommand{\Gal}{\operatorname{Gal}}
\newcommand{\lcm}{\operatorname{lcm}}
\newcommand{\Res}{\operatorname{Res}}

\newtheorem{theorem}{Theorem}[section]
\newtheorem{proposition}[theorem]{Proposition}
\newtheorem{lemma}[theorem]{Lemma}
\newtheorem{corollary}[theorem]{Corollary}
\newtheorem{question}[theorem]{Question}
\newtheorem*{hrsaturation}{Halberstam--Richert saturation theorem}

\theoremstyle{definition}
\newtheorem{definition}[theorem]{Definition}
\newtheorem{example}[theorem]{Example}
\newtheorem{remark}[theorem]{Remark}

\crefname{theorem}{Theorem}{Theorems}
\crefname{proposition}{Proposition}{Propositions}
\crefname{lemma}{Lemma}{Lemmas}
\crefname{corollary}{Corollary}{Corollaries}
\crefname{definition}{Definition}{Definitions}
\crefname{example}{Example}{Examples}
\crefname{remark}{Remark}{Remarks}
\crefname{question}{Question}{Questions}

\title[Universal separation of divisor profiles]{Intersective polynomials
and universal separation of divisor profiles}
\author{Zihan Zhang}

\address{(Zihan Zhang) School of Mathematical Sciences and LPMC, Nankai University, Tianjin
300071, People's Republic of China}
\email{2211056@mail.nankai.edu.cn}

\subjclass[2020]{11A05}
\keywords{divisors, intersective polynomials, almost-prime values}

\begin{document}

\begin{abstract}
We classify universal divisor-profile separation for coprime polynomial pairs
of arbitrary degree and for all pairs of degree at most two.  For
\(A\subset\N\) and \(m\in\Z\), let \(d_A(m)\) count the members of \(A\)
dividing \(m\).  For coprime nonzero \(F,G\in\Z[x]\), universally unbounded
separation between \(d_A(F(n))\) and \(d_A(G(n))\) occurs if and only if one
of \(F,G\) is intersective, that is, has a root modulo every positive integer.
More generally, an intersective factor separated from finitely many polynomial
opponents yields simultaneous one-sided dominance against all of them.

For pairs with common irreducible factors, let \(U,V\) be the products of
the factors occurring only on the two respective sides.  Universal
separation forces \(UV\) to have a root modulo almost every prime;
equivalently, its Galois action has no derangement.  We resolve the remaining
finite \(p\)-adic boundary for all pairs of degree at most two: separation
holds exactly when \(UV\) and at least one of \(F,G\) are intersective, and
the criterion is unchanged by contents or factor multiplicities.  The same
criterion holds, in arbitrary degree, for three linear support factors with
arbitrary positive multiplicities.  The proofs combine uniform almost-prime
values on root progressions with an adaptive local-routing argument.
\end{abstract}

\maketitle

\section{Introduction}

For an infinite set \(A\subset\N=\{1,2,\ldots\}\), S\'ark\H{o}zy asked
whether the difference between the numbers of members of \(A\) dividing
\(n\) and \(n+1\) must be unbounded
\cite{Sarkozy2001}.  Ding answered this question and subsequently classified
all linear pairs \(bn+c,en+f\) with the analogous property
\cite{DingLinear}.  Ding also proved a simultaneous form in which \(n\)
dominates finitely many neighbors \(b_i n\pm1\)
\cite{DingProblem26}.

Relative to these linear results, the present paper classifies coprime pairs
in arbitrary degree, replaces the linear neighbors by finitely many arbitrary
polynomial opponents, isolates the quotient support governing the first
common-factor obstruction, treats three linear support factors with arbitrary
positive multiplicities, and closes all remaining support and multiplicity
patterns in degree at most two.

The modular-root condition now called polynomial intersectivity appears
already in \cite[p.~335, item~(2)]{KamaeMendesFrance}, where the corresponding
polynomial value sets are shown to be van der Corput sets.  Intersectivity
has since served as the exact local condition in polynomial recurrence and
difference-set problems; see
\cite{Lucier2006,BergelsonLeibmanLesigne,Le2010}.
Here it enters a different universal problem.  In the coprime setting it is
the complete threshold for profile separation, whereas in the common-factor
setting its almost-everywhere shadow---prime-covering---first appears as a
necessary Galois condition.  The finite set of primes suppressed by that
condition is precisely where the remaining \(p\)-adic difficulty lies.  To
make this distinction effective in the low-degree cases used below, we also
record a complete finite test for intersectivity through degree three in
\cref{prop:cubic-intersectivity}.

We consider the polynomial extension proposed in
\cite[Problem 3]{DingLinear}.  For any \(A\subset\N\) and \(m\in\Z\), put
\[
  d_A(m)=\#\{a\in A:a\mid m\},
\]
where the cardinality may be infinite.  In particular, \(d_A(0)=|A|\).
When \(A\) is infinite, this requires care if both arguments vanish.  We
therefore use the extended difference
\begin{equation}\label{eq:extended-difference}
  \delta_A(u,v)=
  \begin{cases}
    0,&u=v=0,\\
    \infty,&\text{exactly one of \(u,v\) is zero},\\
    |d_A(u)-d_A(v)|,&uv\ne0.
  \end{cases}
\end{equation}
For \(F,G\in\Z[x]\), let \(\cP(F,G)\) denote the property
\begin{equation}\label{eq:property-P}
  \text{for every infinite \(A\subset\N\)},\qquad
  \limsup_{n\to\infty}\delta_A(F(n),G(n))=\infty.
\end{equation}

The first result gives a complete answer when the polynomials have no common
factor.

\begin{theorem}[Coprime classification]\label{thm:coprime-classification}
Let \(F,G\in\Z[x]\setminus\{0\}\) be coprime in \(\Q[x]\).  Then
\[
  \cP(F,G)
  \quad\Longleftrightarrow\quad
  \text{\(F\) or \(G\) is intersective}.
\]
\end{theorem}

Here a polynomial is \emph{intersective} if it has a root modulo every
positive integer.  The positive half of
\cref{thm:coprime-classification} follows from a stronger simultaneous
statement.

\begin{theorem}[Simultaneous dominance]\label{thm:simultaneous-dominance}
Let \(s\ge1\) and \(F,G_1,\ldots,G_s\in\Z[x]\setminus\{0\}\).  Suppose that
\(F\) has a nonconstant intersective divisor \(H\in\Z[x]\) such that
\[
  \gcd\!\left(H,\prod_{j=1}^sG_j\right)=1
  \qquad\text{in \(\Q[x]\)}.
\]
There is a constant \(C=C(H,G_1,\ldots,G_s)\) with the following
properties.
\begin{enumerate}
  \item If \(A\subset\N\) is finite and \(|A|>C\), then for infinitely many
  positive integers \(n\),
  \[
    d_A(F(n))>\max_{1\le j\le s}d_A(G_j(n)).
  \]
  \item If \(A\subset\N\) is infinite, then
  \[
    \limsup_{n\to\infty}
    \left(d_A(F(n))-\max_{1\le j\le s}d_A(G_j(n))\right)=\infty.
  \]
\end{enumerate}
The integers \(n\) in both assertions may be chosen so that
\(\prod_jG_j(n)\ne0\).
\end{theorem}

Ding's simultaneous theorem \cite[Theorem 4]{DingProblem26} is the model
case with center \(n\) and linear opponents \(b_i n\pm1\).
\Cref{thm:simultaneous-dominance} allows a general intersective center and
arbitrary polynomial opponents, subject to the stated separation condition.

The proof uses a classical saturation theorem of Halberstam and Richert
\cite[Theorem 10.11]{HalberstamRichert}.  Since our root progressions produce
integer-valued polynomials with varying denominators and coefficients, we
give the reduction to the classical integral-polynomial theorem in full.
This also makes clear that the saturation number is uniform in all
coefficients.

When \(F\) and \(G\) share factors, intersectivity may be carried by their
common part.  To state a necessary condition, factor the primitive parts of
\(F\) and \(G\) into primitive irreducibles over \(\Z\).  Let \(C\) be the
product of the irreducible factors occurring on both sides, and let \(U\)
and \(V\) be the products of the factors occurring only in \(F\) and only
in \(G\), respectively; multiplicities are discarded.  Thus, up to signs,
\[
  \radQ(F)=CU,\qquad \radQ(G)=CV,
\]
and \(C,U,V\) are pairwise coprime.

\begin{theorem}[Common-support obstruction]\label{thm:common-support}
Let \(F,G\in\Z[x]\setminus\{0\}\), and let \(C,U,V\) be their support
factors defined above.  If \(UV\) fails to have a root modulo infinitely many
primes, then \(\cP(F,G)\) does not hold.  Consequently,
\[
  \cP(F,G)\quad\Longrightarrow\quad
  UV\text{ has a root modulo all but finitely many primes}.
\]
\end{theorem}

When \(UV\) is nonconstant, let \(L\) be its splitting field, and
let \(\Gamma=\Gal(L/\Q)\) act on its roots.  Away from finitely many primes,
\(UV\) has a root modulo \(p\) exactly when a Frobenius element fixes a root.
Chebotarev therefore yields the following finite-group formulation.

\begin{corollary}[Derangement obstruction]\label[corollary]{cor:derangement}
Suppose that \(UV\) is nonconstant.  If \(\Gamma\) contains an element fixing
no root of \(UV\), then
\(\cP(F,G)\) does not hold.  Thus \(\cP(F,G)\) forces \(\Gamma\) to be the
union of the point stabilizers in its action on the roots of \(UV\).
\end{corollary}

The local--Galois obstruction is not the end of the common-factor theory.
The following exact family lies beyond both sides of the elementary
factor-support sandwich.

\begin{theorem}[Three-linear-support classification]
\label{thm:linear-defects}
Let
\[
  L(x)=\ell x+\lambda,\qquad
  U(x)=ux+\mu,\qquad
  V(x)=vx+\nu
\]
be primitive nonconstant polynomials in \(\Z[x]\), and suppose that they are
pairwise nonassociate.  Let \(\gamma,\delta\in\Z\setminus\{0\}\), let
\(r_L,r_U,s_L,s_V\) be positive integers, and put
\[
  F=\gamma L^{r_L}U^{r_U},\qquad
  G=\delta L^{s_L}V^{s_V}.
\]
Then
\[
  \cP(F,G)\quad\Longleftrightarrow\quad
  \begin{aligned}
    \gcd(|u|,|v|)&=1,\\
    \min\{\gcd(|\ell|,|u|),\gcd(|\ell|,|v|)\}&=1.
  \end{aligned}
\]
Equivalently, \(UV\) and at least one of \(F,G\) are intersective.
\end{theorem}

Thus the criterion depends only on the three squarefree support factors, even
though the four exponents, and hence the degrees, are arbitrary.  The proof
introduces an adaptive local-routing mechanism.  Finitely many resultant
valuations define thresholds at the bad primes, and each prescribed divisor
is routed through the linear branch on which the opposing valuation can be
controlled.  This separates the reusable method from the final degree-two
case analysis.

\begin{corollary}[Complete degree-two classification]
\label[corollary]{cor:degree-two}
Let \(F,G\in\Z[x]\setminus\{0\}\) satisfy
\(\deg F,\deg G\le2\), and let \(C,U,V\) be their support factors.  Then
\[
  \cP(F,G)
  \quad\Longleftrightarrow\quad
  \text{\(UV\) and at least one of \(F,G\) are intersective}.
\]
Moreover,
\[
  \cP(F,G)
  \quad\Longleftrightarrow\quad
  \cP\bigl(\radQ(F),\radQ(G)\bigr).
\]
\end{corollary}

\begin{corollary}[A normalized quadratic family]
\label[corollary]{cor:normalized-linear}
For \(a,b\in\Z\setminus\{0\}\),
\[
  \cP\bigl(x(ax+1),x(bx+1)\bigr)
  \quad\Longleftrightarrow\quad
  \gcd(|a|,|b|)=1\ \text{ and }\ a\ne b.
\]
\end{corollary}

This is the natural exceptional boundary.  Berend and Bilu
\cite{BerendBilu} gave an effective local-and-Galois criterion for a
polynomial to have roots modulo every integer, together with the
corresponding prime-density formula.  More recent classification results for
special families of prime-covering and intersective polynomials appear in
\cite{ElsholtzKlahnTechnau,Mishra2024}.  The point here is that the quotient
support \(UV\), rather than \(F\) or \(G\) separately, governs the first
common-factor obstruction.

The paper is organized as follows.  \Cref{sec:local} fixes the local
language and records elementary obstructions.  \Cref{sec:sieve} proves the
uniform almost-prime and fixed-divisor lemmas.  The simultaneous theorem and
coprime classification are proved in \cref{sec:positive}.  Common supports
and the Galois obstruction are treated in \cref{sec:common}, together with a
low-degree intersectivity test.  The final section proves
\cref{thm:linear-defects,cor:degree-two}, completes the degree-two part of
the exceptional \(p\)-adic boundary, and contrasts a prime-covering but
non-intersective example with a genuinely intersective one.

\section{Local language and elementary obstructions}\label{sec:local}

\begin{definition}
A polynomial \(H\in\Z[x]\) is \emph{intersective} if for every \(m\ge1\)
there exists \(n\in\Z\) such that \(H(n)\equiv0\pmod m\).
\end{definition}

For a nonzero \(P\in\Z[x]\) and \(m\ge1\), write
\[
  Z_P(m)=\{x\bmod m:P(x)\equiv0\pmod m\}.
\]
Intersectivity is equivalent to the existence of a root in \(\Z_p\) for
every prime \(p\).  It is unchanged when polynomial multiplicities are
discarded.

\begin{lemma}[Multiplicity does not affect intersectivity]
\label[lemma]{lem:radical-intersective}
For every nonzero \(P\in\Z[x]\), the polynomial \(P\) is intersective if
and only if its squarefree radical over \(\Q\), represented primitively in
\(\Z[x]\), is intersective.
\end{lemma}

\begin{proof}
Factor the primitive part of \(P\) as
\[
  \prod_{i=1}^r f_i^{e_i}
\]
with distinct primitive irreducibles.  Fix a prime \(p\).  If \(P\) has
roots modulo every power of \(p\), compactness of \(\Z_p\) gives a
\(p\)-adic zero \(z\) of \(P\).  Since \(\Z_p\) is an integral domain,
\(f_i(z)=0\) for some \(i\), and hence the squarefree product
\(\prod_i f_i\) has a \(p\)-adic zero.  Conversely, a zero of the squarefree
product is a zero of \(P\), because the squarefree product divides the
primitive part of \(P\) in \(\Z[x]\).  Applying
this at every prime and using the Chinese remainder theorem proves the
claim.
\end{proof}

The following degree-two criterion is known; see
\cite[Proposition 1 and Appendix B.1]{HamelLyallRice}.  We include the short
proof to specify the normalization used later.

\begin{proposition}[Intersective quadratics]\label[proposition]{prop:quadratic}
Let
\[
  P(x)=ax^2+bx+c\in\Z[x],\qquad a\ne0.
\]
Then \(P\) is intersective if and only if its discriminant is a square and,
on writing its rational roots in lowest terms as
\[
  r_1=\frac{u_1}{v_1},\qquad r_2=\frac{u_2}{v_2},
  \qquad v_1,v_2>0,
\]
one has \(\gcd(v_1,v_2)=1\).
\end{proposition}

\begin{proof}
If the discriminant is not a square, its associated nontrivial quadratic
character is negative at infinitely many primes.  At every prime away from
\(2a\) and the discriminant where that character is negative, \(P\) has no
root.  Thus an intersective quadratic factors over \(\Q\).

By Gauss's lemma, write
\[
  P(x)=K(A_1x+B_1)(A_2x+B_2)
\]
with \(K\in\Z\setminus\{0\}\) and
\(\gcd(A_i,B_i)=1\).  If a prime \(p\) divides both \(A_1\) and \(A_2\),
then both linear factors are \(p\)-adic units for every integer \(x\).
Consequently \(v_p(P(x))=v_p(K)\) for all \(x\), so \(P\) has no root
modulo \(p^{v_p(K)+1}\).

Conversely, suppose \(\gcd(A_1,A_2)=1\).  For every prime power \(p^k\),
at least one \(A_i\) is a unit modulo \(p\), and the corresponding linear
factor has a root modulo \(p^k\).  The Chinese remainder theorem combines
these choices for an arbitrary modulus.  Finally, the reduced root
denominators are \(|A_1|\) and \(|A_2|\).
\end{proof}

\begin{proposition}[Two missing local roots]
\label[proposition]{prop:both-nonintersective}
Let \(F,G\in\Z[x]\setminus\{0\}\).  If neither polynomial is intersective,
then \(\cP(F,G)\) does not hold.
\end{proposition}

\begin{proof}
Choose \(q,r\ge2\) such that \(F\) has no root modulo \(q\) and \(G\) has
no root modulo \(r\).  Put \(M=\lcm(q,r)\) and
\[
  A=\{M^k:k\ge1\}.
\]
No member of \(A\) divides a value of \(F\) or \(G\).  The same local
obstructions rule out integer zeros.  Hence
\(\delta_A(F(n),G(n))=0\) for every \(n\).
\end{proof}

\begin{proposition}[Identical local root sets]
\label[proposition]{prop:identical-roots}
Suppose that \(F,G\in\Z[x]\setminus\{0\}\) and there is an infinite
\(A\subset\N\) such that
\[
  Z_F(a)=Z_G(a)\qquad(a\in A).
\]
Then \(\cP(F,G)\) does not hold.
\end{proposition}

\begin{proof}
For every \(a\in A\) and \(n\in\Z\),
\[
  a\mid F(n)\quad\Longleftrightarrow\quad a\mid G(n).
\]
Thus the two finite divisor counts agree whenever \(F(n)G(n)\ne0\).
If one value is zero, every \(a\in A\) divides the other value; since
\(A\) is unbounded, the other value is zero as well.  The definition
\eqref{eq:extended-difference} now gives
\(\delta_A(F(n),G(n))=0\) for every \(n\).
\end{proof}

\begin{remark}[Agreement and disagreement of local root sets]
\label[remark]{rem:root-set-levels}
Distinct polynomials can have identical root sets along an infinite family of
moduli without having identical root data at every modulus.  For example, let
\(F(x)=x\) and \(G(x)=x^2\).  For every prime \(p\),
\[
  Z_F(p)=Z_G(p)=\{0\bmod p\},
\]
so \cref{prop:identical-roots} applies with \(A\) equal to the set of primes.
On the other hand, for \(k\ge2\),
\[
  Z_F(p^k)=\{0\bmod p^k\},\qquad
  Z_G(p^k)=
  \{x\bmod p^k:p^{\lceil k/2\rceil}\mid x\},
\]
and these sets are different.  Thus the hypothesis of
\cref{prop:identical-roots} is a condition on the selected moduli; it neither
requires the polynomials to be associates nor asserts equality at all
prime-power levels.
\end{remark}

\section{Uniform almost-prime values on root progressions}
\label{sec:sieve}

For a nonzero integer \(m\), let \(\Omega(m)\) be the number of prime
factors of \(|m|\), counted with multiplicity.  If
\(Q\in\Q[x]\) is integer-valued, define its fixed divisor by
\[
  \fd(Q)=\gcd\{Q(n):n\in\Z\},
\]
taking the positive gcd.

We use the following explicit form of the classical saturation theorem.  It
is the first of the two admissible numerical alternatives in
\cite[Theorem~10.11, pp.~310--311]{HalberstamRichert}.

\begin{hrsaturation}
Let \(f_1,\ldots,f_g\in\Z[x]\) be distinct irreducible polynomials with
positive leading coefficients.  Put
\[
  f=\prod_{i=1}^g f_i,\qquad E=\deg f,
  \qquad \rho(p)=\#\{a\bmod p:f(a)\equiv0\pmod p\}.
\]
Assume that \(\rho(p)<p\) for every prime \(p\).  If the integer \(b\)
satisfies
\begin{equation}\label{eq:HR-bound}
  b>E-1+g\sum_{j=1}^g\frac1j
      +g\log\left(\frac{2E}{g}+\frac1{g+1}\right),
\end{equation}
then there is a constant \(\delta=\delta(b,f)>0\) such that, as
\(x\to\infty\),
\[
  \#\{1\le n\le x:f(n)=P_b\}
  \ge \delta\frac{x}{(\log x)^g}
      \left(1+O_f\left(\frac1{\log\log x}\right)\right).
\]
Here \(P_b\) denotes a positive integer with at most \(b\) prime factors,
counted with multiplicity.  In particular, \(\Omega(f(n))\le b\) for
infinitely many positive integers \(n\).
\end{hrsaturation}

The threshold in \eqref{eq:HR-bound} depends only on the number and degrees
of the irreducible factors.  The density constant and the starting point may
depend on their coefficients; only the almost-prime bound will need to be
uniform below.

\begin{lemma}[Uniform saturation for integer-valued polynomials]
\label[lemma]{lem:uniform-saturation}
For each \(D\ge1\), there is an integer \(R_D\) with the following property.
If \(Q\in\Q[x]\) is a nonconstant integer-valued polynomial of degree at
most \(D\) and \(\fd(Q)=1\), then there are infinitely many positive
integers \(y\) such that
\[
  Q(y)\ne0,\qquad \Omega(Q(y))\le R_D.
\]
One possible choice is
\[
  R_D=Db_D,\qquad
  b_D=1+\left\lfloor
  D-1+D\sum_{j=1}^D\frac1j+D\log(2D+1)
  \right\rfloor.
\]
\end{lemma}

\begin{proof}
We reduce the assertion to the preceding saturation theorem.  Let
\(d=\deg Q\le D\).
Choose a positive integer \(c\) such that
\[
  P(x)=cQ(x)\in\Z[x].
\]
Since \(\fd(Q)=1\), one has \(\fd(P)=c\).  For each prime \(p\mid c\),
choose \(a_p\in\Z\) such that
\[
  v_p(P(a_p))=v_p(c).
\]
By the Chinese remainder theorem, choose \(a\) satisfying
\[
  a\equiv a_p\pmod {p^{v_p(c)+1}}\qquad(p\mid c).
\]
We claim that
\[
  q(z)=Q(a+cz)=\frac{P(a+cz)}{c}
\]
belongs to \(\Z[z]\) and has no fixed prime divisor.

The difference \(P(a+cz)-P(a)\) has every coefficient divisible by \(c\),
while \(P(a)/c=Q(a)\) is an integer; hence \(q\in\Z[z]\).
If \(p\mid c\), congruence preservation and the choice of \(a\) give
\[
  q(0)=\frac{P(a)}c\not\equiv0\pmod p.
\]
If \(p\nmid c\) divided every value of \(q\), then, since multiplication by
\(c\) permutes the residue classes modulo \(p\), the prime \(p\) would divide
\(P(n)\) for every integer \(n\).  This contradicts \(\fd(P)=c\).  The
claim follows.

Factor
\[
  q(z)=\pm\prod_{i=1}^r f_i(z)^{e_i},
\]
where the \(f_i\) are distinct primitive irreducible polynomials with
positive leading coefficients.  There is no nontrivial constant content,
because \(q\) has no fixed prime divisor.  The squarefree product
\[
  q_{\mathrm{sf}}(z)=\prod_{i=1}^r f_i(z)
\]
also has no fixed prime divisor.

For a prime \(p\), let \(\rho(p)\) denote the number of roots of
\(q_{\mathrm{sf}}\) modulo \(p\).  The absence of a fixed prime divisor says
\(\rho(p)<p\) for every \(p\).  For each admissible integer \(b\), the
saturation theorem supplies infinitely many positive \(z\) for which
\[
  \Omega(q_{\mathrm{sf}}(z))\le b.
\]
Its positivity convention causes no issue, since
\(q_{\mathrm{sf}}(z)>0\) for all sufficiently large positive \(z\).  Since
\(\sum_i e_i\deg f_i=d\), we have \(e_i\le d\), and
therefore
\[
  \Omega(q(z))
  =\sum_{i=1}^r e_i\Omega(f_i(z))
  \le d\,\Omega(q_{\mathrm{sf}}(z)).
\]
Let \(E=\deg q_{\mathrm{sf}}\le D\) and \(r\le D\).  The theorem allows any
integer strictly larger than
\[
  E-1+r\sum_{j=1}^r\frac1j+
  r\log\left(\frac{2E}{r}+\frac1{r+1}\right).
\]
This quantity is strictly smaller than \(b_D\), so \(b=b_D\) is admissible.
Consequently,
\[
  \Omega(q(z))\le d b_D\le D b_D=R_D.
\]
The notation \(P_r\) in
\cite[p. 242]{HalberstamRichert} counts prime factors with multiplicity,
which is the function \(\Omega\) used here.  Finally, the corresponding
arguments \(y=a+cz\) are distinct and positive for all sufficiently large
positive \(z\), completing the proof.
\end{proof}

\begin{remark}
The starting point from which the infinitely many \(y\) occur may depend on
the coefficients of \(Q\).  Only the bound \(R_D\) is uniform.  This is
exactly the uniformity required below, because the progression polynomial
changes with the number of forced divisors.
\end{remark}

\begin{lemma}[Fixed divisors on root progressions]
\label[lemma]{lem:fixed-divisor}
Let \(H,K\in\Z[x]\setminus\{0\}\) be coprime in \(\Q[x]\).  There is a
constant \(C_0=C_0(H,K)\) such that, for all \(a\in\Z\) and \(M\ge1\)
satisfying \(M\mid H(a)\),
\[
  \fd\bigl(K(a+M\,\cdot)\bigr)\le C_0.
\]
\end{lemma}

\begin{proof}
There are \(B,C\in\Z[x]\) and a nonzero integer \(R\) such that
\begin{equation}\label{eq:bezout}
  B(x)H(x)+C(x)K(x)=R.
\end{equation}
Fix \(a,M\) with \(M\mid H(a)\), and let
\[
  \Delta=\fd\bigl(K(a+M\,\cdot)\bigr).
\]
For a prime \(p\), write
\[
  \delta_p=v_p(\Delta)
  =\min_{y\in\Z}v_p(K(a+My)).
\]

If \(p\nmid M\), the progression \(a+My\) covers every residue class
modulo every power of \(p\).  Hence
\[
  \delta_p\le v_p(\fd(K)).
\]

Suppose \(p\mid M\), and set \(m=v_p(M)\).  Congruence preservation gives
\(p^m\mid H(a+My)\) for every \(y\).  From \eqref{eq:bezout},
\[
  p^{\min(m,\delta_p)}\mid R.
\]
If \(m>v_p(R)\), this gives \(\delta_p\le v_p(R)\).  If
\(1\le m\le v_p(R)\), then \(p\mid R\) and there are only finitely many
possible pairs consisting of \(m\) and a residue class
\(a\bmod p^m\) with \(H(a)\equiv0\pmod {p^m}\).  For each such pair,
\[
  e_{p,m,a}:=\min_{x\equiv a\;(\mathrm{mod}\;p^m)}v_p(K(x))
\]
is finite, since \(K\) is not the zero polynomial.  Moreover this minimum
equals \(\delta_p\).  Indeed, write \(M=p^mM'\) with \(p\nmid M'\), choose
an \(x\) attaining the displayed minimum, and use that \(M'\) is a
\(p\)-adic unit to choose \(y\) with
\[
  a+My\equiv x\pmod {p^{e_{p,m,a}+1}}.
\]
Then \(v_p(K(a+My))=e_{p,m,a}\).  Taking the maximum over these finitely
many cases bounds \(\delta_p\).  Only primes dividing
\(R\,\fd(K)\) can contribute, so their bounds combine to a constant
\(C_0(H,K)\).
\end{proof}

\section{Simultaneous dominance and the coprime classification}
\label{sec:positive}

\begin{proof}[Proof of \cref{thm:simultaneous-dominance}]
Put
\[
  K(x)=\prod_{j=1}^sG_j(x).
\]
First suppose that \(K\) is nonconstant.  Let
\[
  A_t=\{a_1,\ldots,a_t\}\subset\N,\qquad
  M_t=\prod_{i=1}^t a_i.
\]
Here \(A_t=A\) in the finite case, while in the infinite case the elements
of \(A\) are listed increasingly and \(A_t\) consists of the first \(t\).

Choose \(u_t\in\Z\) such that \(M_t\mid H(u_t)\).  Then for every
\(y\in\Z\),
\[
  M_t\mid H(u_t+M_ty)\mid F(u_t+M_ty).
\]
In particular, every member of \(A_t\) divides the corresponding value of
\(F\).

Let
\[
  \Delta_t=\fd\bigl(K(u_t+M_t\,\cdot)\bigr).
\]
By \cref{lem:fixed-divisor}, \(1\le\Delta_t\le C_0(H,K)\).  The polynomial
\[
  Q_t(y)=\frac{K(u_t+M_ty)}{\Delta_t}
\]
is nonconstant, integer-valued, has degree \(\deg K\), and has fixed divisor
one.  By \cref{lem:uniform-saturation}, there are arbitrarily large positive
\(y\) such that
\[
  Q_t(y)\ne0,\qquad \Omega(Q_t(y))\le R_{\deg K}.
\]
For \(n=u_t+M_ty\), all \(G_j(n)\) are nonzero.  Moreover,
\[
  \Omega(K(n))
  \le R_{\deg K}+
  \max_{1\le m\le C_0(H,K)}\Omega(m)=:B,
\]
and hence, for every \(j\),
\begin{equation}\label{eq:opponent-bound}
  d_A(G_j(n))
  \le\tau(|G_j(n)|)
  \le\tau(|K(n)|)
  \le2^B=:C.
\end{equation}
The constant is independent of \(t\) and \(A\).

If \(A\) is finite and \(|A|>C\), use \(A_t=A\) to obtain
\[
  d_A(F(n))\ge|A|>C\ge\max_jd_A(G_j(n))
\]
for infinitely many \(n\).  If \(A\) is infinite, choose the saturating
value of \(y\) large enough at each stage that \(n\to\infty\).  Then
\[
  d_A(F(n))-\max_jd_A(G_j(n))\ge t-C,
\]
which proves the second assertion.

If \(K\) is a nonzero constant, the same proof stops before the sieve step:
all opponent divisor counts are bounded by the divisor counts of their fixed
nonzero values.
\end{proof}

\begin{proof}[Proof of \cref{thm:coprime-classification}]
If neither \(F\) nor \(G\) is intersective, the property fails by
\cref{prop:both-nonintersective}.  Conversely, suppose for example that
\(F\) is intersective.  It is nonconstant, and
\cref{thm:simultaneous-dominance} applies with \(H=F\), \(s=1\), and
opponent \(G\).  It gives an unbounded positive signed difference, hence
\(\cP(F,G)\).  The case in which \(G\) is intersective is symmetric.
\end{proof}

\begin{corollary}[Two-sided oscillation]\label[corollary]{cor:oscillation}
Let \(F,G\in\Z[x]\setminus\{0\}\) be coprime in \(\Q[x]\), and suppose both
are intersective.  For every infinite \(A\subset\N\),
\[
  \limsup_{n\to\infty}\bigl(d_A(F(n))-d_A(G(n))\bigr)=\infty
\]
and
\[
  \liminf_{n\to\infty}\bigl(d_A(F(n))-d_A(G(n))\bigr)=-\infty,
\]
where the two sequences may be restricted to integers at which both
polynomial values are nonzero.
\end{corollary}

\begin{proof}
Apply \cref{thm:simultaneous-dominance} first to \(F\) against \(G\), and
then to \(G\) against \(F\).
\end{proof}

\begin{corollary}[Zero and constant cases]
\label[corollary]{cor:zero-constant}
Let \(F,G\in\Z[x]\).
\begin{enumerate}
  \item If exactly one of \(F,G\) is the zero polynomial, then
  \(\cP(F,G)\) holds.
  \item If both are zero, then \(\cP(F,G)\) fails.
  \item If \(C\ne0\) is constant and \(P\) is nonzero and nonconstant, then
  \(\cP(C,P)\) holds if and only if \(P\) is intersective.
\end{enumerate}
\end{corollary}

\begin{proof}
The first two assertions follow immediately from
\eqref{eq:extended-difference}, since a nonzero polynomial has only finitely
many integer zeros.  The third assertion is the coprime classification.
\end{proof}

\section{Recognizing intersectivity and common-support obstructions}
\label{sec:common}

The definition of intersectivity is elementary, but its effective
recognition has two distinct layers.  Replace a polynomial by its primitive
squarefree radical and factor it over \(\Q\) as
\[
  R=h_1\cdots h_r.
\]
If \(L\) is the splitting field, \(\Gamma=\Gal(L/\Q)\), and \(H_i\) is the
stabilizer of a root of \(h_i\), then the good-prime condition is the finite
permutation-group covering
\[
  \Gamma=\bigcup_{i=1}^r\ \bigcup_{\sigma\in\Gamma}
  \sigma H_i\sigma^{-1}.
\]
Berend and Bilu \cite{BerendBilu} proved that this condition, together with
solvability at one explicitly computable modulus supported on the
resultants \(\Res(h_i,h_i')\), is equivalent to intersectivity.  Thus the
problem is decidable in principle, but the calculation can require
factorization over \(\Q\), a splitting field and its permutation group, and
separate \(p\)-adic analysis at the finitely many bad primes.

\begin{definition}
A nonzero polynomial \(R\in\Z[x]\) is \emph{prime-covering} if it has a
root modulo all but finitely many primes.
\end{definition}

Every intersective polynomial is prime-covering.  For a squarefree
polynomial, the converse holds at every good prime: a root modulo \(p\) is
simple and hence lifts to a root in \(\Z_p\).  The gap between the two notions
is therefore concentrated at finitely many primes, but it is genuinely
\(p\)-adic: roots must persist modulo \(p^k\) for every \(k\), not merely
modulo \(p\).  Equivalently, prime-covering is the Galois layer of the
Berend--Bilu criterion, whereas intersectivity also includes its finite local
layer.  The distinction is illustrated concretely in
\cref{ex:236}.  Recent work gives more explicit criteria and constructions
for special factorization patterns
\cite{ElsholtzKlahnTechnau,MishraQuadratic,MishraMinimal,Mishra2024}, but
there is no comparably short coefficient criterion in unrestricted degree.

The preceding framework becomes completely elementary through degree three.

\begin{proposition}[Complete intersectivity test through degree three]
\label[proposition]{prop:cubic-intersectivity}
Let \(P\in\Z[x]\setminus\{0\}\) have degree at most three, and let \(R\)
be its squarefree radical over \(\Q\), represented primitively in \(\Z[x]\).
Then \(P\) is intersective if and only if one of the following alternatives
holds.
\begin{enumerate}
  \item The polynomial \(R\) splits over \(\Q\).  On writing
  \[
    R(x)=\pm\prod_{i=1}^r(a_i x+b_i),
    \qquad 1\le r\le3,
  \]
  with each factor primitive, one has
  \[
    \gcd(|a_1|,\ldots,|a_r|)=1.
  \]
  \item The polynomial \(R\) has degree three and
  \[
    R(x)=\pm L(x)Q(x),\qquad L(x)=ax+b,
  \]
  where \(L\) is primitive and \(Q\in\Z[x]\) is primitive and irreducible
  quadratic.  For every prime \(p\mid a\), the polynomial \(Q\) has a root
  in \(\Z_p\).
\end{enumerate}
The test in \textup{(ii)} is finite.  If
\[
  s_p=v_p\bigl(\lvert\Res(Q,Q')\rvert\bigr),
\]
then its local condition is equivalent to the solvability of
\[
  Q(x)\equiv0\pmod {p^{2s_p+1}}
\]
for every prime \(p\mid a\).
\end{proposition}

\begin{proof}
By \cref{lem:radical-intersective}, it suffices to consider \(R\).  A
primitive linear polynomial \(ax+b\) has a root in \(\Z_p\) exactly when
\(p\nmid a\).  Since \(\Z_p\) is an integral domain, a product of primitive
linear factors has a root in \(\Z_p\) exactly when at least one of their
leading coefficients is a \(p\)-adic unit.  Requiring this for every prime
is precisely the gcd condition in \textup{(i)}.

Suppose next that \(R=LQ\) as in \textup{(ii)}.  If \(p\nmid a\), then
\(L\) has a root in \(\Z_p\).  If \(p\mid a\), primitivity gives
\(p\nmid b\), so \(L(x)\) is a \(p\)-adic unit for every
\(x\in\Z_p\); at such a prime, \(R\) has a \(p\)-adic root exactly when
\(Q\) does.  This proves the stated local criterion.

To verify the finite congruence test, put
\(S=\Res(Q,Q')\ne0\).  There are \(A,B\in\Z[x]\) such that
\[
  A(x)Q(x)+B(x)Q'(x)=S.
\]
If \(Q(x)\equiv0\pmod {p^{2s_p+1}}\), the identity implies
\(v_p(Q'(x))\le s_p\).  Hence
\[
  v_p(Q(x))>2v_p(Q'(x)),
\]
and the strong form of Hensel's lemma gives a zero of \(Q\) in \(\Z_p\).
The converse follows by reducing a \(p\)-adic zero modulo
\(p^{2s_p+1}\).

It remains to exclude the other factorizations.  A squarefree polynomial of
degree at most three with no linear factor is irreducible quadratic or
irreducible cubic.  Its nontrivial transitive Galois action contains a
derangement, so Chebotarev gives infinitely many primes modulo which it has
no root.  It is not prime-covering and therefore not intersective.  The two
listed alternatives exhaust all remaining factorizations.
\end{proof}

\begin{example}[A cubic intersective polynomial without an integral root]
\label[example]{ex:cubic-intersective}
The polynomial
\[
  P(x)=(2x+1)(x^2+x+2)
\]
has no integral root.  Its quadratic factor is irreducible over \(\Q\), and
the only prime dividing the leading coefficient of \(2x+1\) is \(2\).
Modulo \(2\), the quadratic has the simple root \(0\), so Hensel's lemma
gives a root in \(\Z_2\).  \Cref{prop:cubic-intersectivity}\textup{(ii)}
therefore shows that \(P\) is intersective.
\end{example}

The proposition is a complete test for the intersectivity of a single
polynomial through degree three.  It is not a degree-three classification of
\(\cP(F,G)\); common-support pairs in that range already meet the higher
local-routing difficulty described in \cref{q:exceptional-boundary}.

We make the factor-support normalization precise.  For a nonzero
\(P\in\Z[x]\), write
\[
  P=c_P\prod_{\phi\in\cI_P}\phi^{e_P(\phi)},
\]
where \(c_P\in\Z\setminus\{0\}\) and the \(\phi\) are pairwise
nonassociate primitive irreducible polynomials with positive leading
coefficients.  For \(F,G\ne0\), set
\[
  \cI_C=\cI_F\cap\cI_G,\qquad
  \cI_U=\cI_F\setminus\cI_G,\qquad
  \cI_V=\cI_G\setminus\cI_F,
\]
and define
\[
  C=\prod_{\phi\in\cI_C}\phi,\qquad
  U=\prod_{\phi\in\cI_U}\phi,\qquad
  V=\prod_{\phi\in\cI_V}\phi.
\]
An empty product is \(1\).  Then \(C,U,V\) are pairwise coprime and
\[
  \radQ(F)=CU,\qquad \radQ(G)=CV
\]
up to multiplication by \(-1\).

For a nonconstant squarefree polynomial, prime-covering is a Galois
condition.  In the monic integral setting, a prime-covering polynomial with
no integer root is called \emph{exceptional}; see
\cite{ElsholtzKlahnTechnau}.

\begin{proof}[Proof of \cref{thm:common-support}]
Suppose that \(UV\) has no root modulo any prime in an infinite set
\(S\).  Delete from \(S\) the finitely many primes dividing \(c_Fc_G\).
For \(p\in S\), neither \(U\) nor \(V\) has a root modulo \(p\).  Therefore,
for every integer \(n\),
\[
  p\mid F(n)
  \quad\Longleftrightarrow\quad
  p\mid C(n)
  \quad\Longleftrightarrow\quad
  p\mid G(n).
\]
Taking \(A=S\), viewed as an infinite set of positive integers, gives
\[
  Z_F(p)=Z_G(p)\qquad(p\in A).
\]
\Cref{prop:identical-roots} proves the result.
\end{proof}

We recall the precise Chebotarev--permutation correspondence used next.  For a
finite Galois extension \(K/\Q\), the Chebotarev density theorem assigns to
each conjugacy class \(\mathcal C\subset\Gal(K/\Q)\) a set of unramified
primes of natural density \(|\mathcal C|/|\Gal(K/\Q)|\); see
\cite[Section~2, Chebotarev theorem]{ElsholtzKlahnTechnau}.  If \(K\) is the
splitting field of a squarefree polynomial \(R\), then, outside the primes
dividing its leading coefficient or discriminant, the cycle type of Frobenius
is the factorization type of \(R\bmod p\).  Hence \(R\bmod p\) has a linear
factor exactly when Frobenius fixes a root.  An element of this permutation
action with no fixed root is a \emph{derangement}; thus the primes at which
\(R\) has no root are exactly the primes belonging to derangement conjugacy
classes, apart from the finite exceptional set.  The resulting fixed-point
density formula is stated explicitly in
\cite[Theorem~2, p.~1666]{BerendBilu} and
\cite[Theorem~2.1]{ElsholtzKlahnTechnau}.

\begin{proposition}[Frobenius fixed-point density]
\label[proposition]{prop:frobenius}
Let \(R\in\Z[x]\) be nonconstant and squarefree, let \(L\) be its splitting
field, and let \(\Gamma=\Gal(L/\Q)\) act on the roots of \(R\).  The set of
primes modulo which \(R\) has a root has natural density
\[
  \frac{\#\{\sigma\in\Gamma:\sigma\text{ fixes a root of }R\}}
       {|\Gamma|}.
\]
\end{proposition}

\begin{proof}
The preceding correspondence identifies the required primes with the union
of the conjugacy classes whose elements fix a root.  Summing their
Chebotarev densities gives the displayed quotient; omitting finitely many bad
primes does not change the density.
\end{proof}

\begin{proof}[Proof of \cref{cor:derangement}]
If \(\sigma\in\Gamma\) fixes no root, then the same is true throughout its
conjugacy class.  By \cref{prop:frobenius}, a positive-density set of primes
has no root of \(UV\).  Apply \cref{thm:common-support}.
\end{proof}

\begin{corollary}[A factor-support sandwich]
\label[corollary]{cor:sandwich}
Let \(C,U,V\) be the support factors of \(F,G\ne0\).
\begin{enumerate}
  \item If \(U\) is nonconstant and intersective, or if \(V\) is
  nonconstant and intersective, then \(\cP(F,G)\) holds.
  \item If \(\cP(F,G)\) holds, then \(UV\) is prime-covering; equivalently,
  when \(UV\) is nonconstant, its Galois action has no derangement.
\end{enumerate}
\end{corollary}

\begin{proof}
The polynomial \(U\) divides \(F\) in \(\Z[x]\) and is coprime to \(G\);
the analogous statement holds for \(V\).  Thus the first assertion follows
from \cref{thm:simultaneous-dominance}.  The second is
\cref{thm:common-support,cor:derangement}.  For the stated equivalence,
outside the finite exceptional set in the Frobenius description, every
Frobenius element fixes a root when the Galois action has no derangement;
hence every such prime has a root of \(UV\).
\end{proof}

When \(F\) and \(G\) are coprime, the common factor \(C\) is \(1\).
In that case \cref{prop:both-nonintersective} closes the gap between the two
sides of \cref{cor:sandwich}, giving
\cref{thm:coprime-classification}.  For noncoprime pairs, the remaining case
is genuinely different: the common support can make both original
polynomials intersective even when neither unique factor \(U\) nor \(V\) is
intersective.

\section{The exceptional p-adic boundary}
\label{sec:boundary}

Prime-covering is weaker than intersectivity.  The distinction is already
visible in a classical elementary example.

\begin{example}[Prime-covering but not intersective]\label[example]{ex:236}
Let
\[
  E(x)=(x^2-2)(x^2-3)(x^2-6).
\]
For every odd prime \(p\), if neither \(2\) nor \(3\) is a quadratic residue
modulo \(p\), then \(6\) is a quadratic residue.  At \(p=2\) or \(p=3\),
the residue \(x=0\) is a root, so \(E\) has a root modulo every prime.

On the other hand, if \(x\) is even then
\[
  E(x)\equiv4\pmod8,
\]
while if \(x\) is odd then
\[
  E(x)\equiv6\pmod8.
\]
Thus \(E\) has no root modulo \(8\) and is not intersective.  This example
is discussed in \cite{ElsholtzKlahnTechnau}.
\end{example}

A genuine intersective polynomial without a rational root is
\[
  (x^2-13)(x^2-17)(x^2-221);
\]
see \cite[Example 1]{BerendBilu}.  These two examples show why the
derangement condition alone cannot settle the common-factor case: it sees
almost all primes, whereas \(\cP(F,G)\) may also be controlled by a finite
set of bad primes and their full \(p\)-adic root profiles.

We now prove the common-factor classification stated in the introduction.

\begin{proof}[Proof of \cref{thm:linear-defects}]
If a prime \(p\) divides both \(u\) and \(v\), then \(U(n)\) and \(V(n)\)
are \(p\)-adic units for every \(n\), by primitivity.  Use the lacunary test
set
\[
  \mathcal A=\{p^{3^j}:j\ge1\},
\]
and write
\[
  N(X)=\#\{j\ge1:3^j\le X\}\qquad(X\ge0).
\]
For all sufficiently large \(n\), the polynomial values are nonzero.  With
\(\tau=v_p(L(n))\), the two divisor counts are then
\[
  N\bigl(v_p(\gamma)+r_L\tau\bigr)
  \quad\text{and}\quad
  N\bigl(v_p(\delta)+s_L\tau\bigr).
\]
Their difference is uniformly bounded in \(\tau\), since
\(N(X)=\log_3(X+1)+O(1)\) and the ratio of the two affine arguments after
adding one is bounded above and below.  Thus \(\cP(F,G)\) fails.

By \cref{lem:radical-intersective,prop:quadratic},
\[
  F\text{ is intersective}\quad\Longleftrightarrow\quad
  \gcd(|\ell|,|u|)=1,
\]
and
\[
  G\text{ is intersective}\quad\Longleftrightarrow\quad
  \gcd(|\ell|,|v|)=1.
\]
If both gcds exceed one, the property fails by
\cref{prop:both-nonintersective}.  This proves the necessity of the stated
conditions.

For sufficiency, suppose by symmetry that
\[
  \gcd(|\ell|,|u|)=\gcd(|u|,|v|)=1.
\]
Put
\[
  R=\Res(L,U)=\ell\mu-u\lambda\ne0,
  \qquad h_q=v_q(R),\qquad
  \beta_q=v_q(\delta)\quad(q\mid v).
\]
For an arbitrary infinite \(\mathcal A\subset\N\), define
\[
  \mathcal A_{\le}=
  \{m\in\mathcal A:
    v_q(m)\le\beta_q+s_Lh_q\text{ for every prime }q\mid v\}.
\]
When \(v=\pm1\), the condition is empty and
\(\mathcal A_{\le}=\mathcal A\).

Suppose first that \(\mathcal A_{\le}\) is infinite.  Choose distinct
\(m_1,\ldots,m_t\) in this set, put
\[
  r_i=\frac{m_i}{\gcd(m_i,|\delta|)},
\]
and write
\[
  d_i=\prod_{q\mid v}q^{v_q(r_i)},\qquad
  c_i=\frac{r_i}{d_i}.
\]
Put
\[
  D=\lcm(d_1,\ldots,d_t),\qquad
  C=\lcm(c_1,\ldots,c_t),\qquad
  B=\prod_{q\mid v}q^{s_Lh_q}.
\]
Thus \(D\mid B\), \(\gcd(D,C)=1\), and \(\gcd(v,C)=1\).
Also \(\gcd(\ell,D)=1\).  Indeed, if a prime \(q\) divided both
\(\ell\) and \(D\), then \(q\mid v\) and \(h_q\ge1\).  But
\(q\nmid u\) and \(q\nmid\lambda\), so
\[
  R=\ell\mu-u\lambda\not\equiv0\pmod q,
\]
a contradiction.

The Chinese remainder theorem now gives \(n_0\in\Z\) such that
\[
  L(n_0)\equiv0\pmod D,\qquad
  V(n_0)\equiv0\pmod C.
\]
Set \(M=DC\).  Since \(r_i\mid L(n_0+My)V(n_0+My)\), every \(m_i\) divides
\(G(n_0+My)\), for every integer \(y\).  Moreover,
\[
  M\mid D\,V(n_0).
\]
The polynomial \(H_D=D V\) is coprime to
\(F=\gamma L^{r_L}U^{r_U}\).  Since \(D\mid B\), only finitely many
\(H_D\) occur.  Taking the maximum in
\cref{lem:fixed-divisor}, there is a constant \(C_*\) such that
\[
  \Delta:=\fd\bigl(F(n_0+M\,\cdot)\bigr)\le C_*.
\]
The polynomial \(F(n_0+My)/\Delta\) is integer-valued, has degree
\(r_L+r_U\), and has fixed divisor one.  By
\cref{lem:uniform-saturation}, there are
arbitrarily large positive \(y\), avoiding the finitely many roots of \(G\),
such that
\[
  \begin{aligned}
    F(n_0+My)G(n_0+My)&\ne0,\\
    \Omega(F(n_0+My))
      &\le R_{r_L+r_U}+\max_{1\le d\le C_*}\Omega(d)=:B_0,
  \end{aligned}
\]
where \(B_0\) is independent of \(t\) and the selected elements.  Hence
\[
  d_{\mathcal A}(G(n_0+My))-d_{\mathcal A}(F(n_0+My))
  \ge t-2^{B_0}.
\]
Choosing these arguments increasingly with \(t\) proves the property in this
case.

Suppose now that \(\mathcal A_{\le}\) is finite.  Choose distinct
\(m_1,\ldots,m_t\in\mathcal A\setminus\mathcal A_{\le}\), write
\[
  d_i=\prod_{p\mid u}p^{v_p(m_i)},\qquad
  c_i=\frac{m_i}{d_i},
\]
and put
\[
  D=\lcm(d_1,\ldots,d_t),\qquad
  C=\lcm\left(c_1,\ldots,c_t,
      \prod_{q\mid v}q^{h_q+1}\right).
\]
The assumptions give
\[
  \gcd(\ell,D)=\gcd(u,C)=\gcd(D,C)=1.
\]
Choose \(n_0\) such that
\[
  L(n_0)\equiv0\pmod D,\qquad
  U(n_0)\equiv0\pmod C,
\]
and set \(M=DC\).  Along the progression \(n=n_0+My\), every \(m_i\)
divides \(F(n)\).

Fix a prime \(q\mid v\).  Since \(q^{h_q+1}\mid U(n)\), the identity
\[
  uL(x)-\ell U(x)=-R
\]
and \(q\nmid u\) give
\[
  v_q(L(n))=h_q.
\]
Also \(V(n)\equiv\nu\not\equiv0\pmod q\), because \(V\) is primitive.
Consequently,
\[
  v_q(G(n))=\beta_q+s_Lh_q\qquad(q\mid v).
\]
Every member of \(\mathcal A\setminus\mathcal A_{\le}\) has
\(q\)-adic valuation greater than \(\beta_q+s_Lh_q\) for at least one
\(q\mid v\).
None can divide \(G(n)\).  Taking positive representatives
\(n_t\to\infty\), avoiding the finitely many roots of \(FG\), yields
\[
  d_{\mathcal A}(F(n_t))-d_{\mathcal A}(G(n_t))
  \ge t-|\mathcal A_{\le}|,
\]
which completes the proof.
\end{proof}

\begin{proof}[Proof of \cref{cor:degree-two}]
If \(C=1\), the polynomials are coprime.  By
\cref{thm:coprime-classification}, the property holds exactly when one of
\(F,G\) is intersective; in that event their product \(UV\) is automatically
intersective.

Suppose that \(C\) is nonconstant.  If \(U=V=1\), take for
\(\mathcal A\) all primes not dividing the two polynomial contents.  For
every \(p\in\mathcal A\), divisibility by \(p\) depends only on the common
squarefree support, so
\[
  Z_F(p)=Z_G(p).
\]
The property fails by \cref{prop:identical-roots}, while \(UV=1\) is not
intersective.

Suppose next that \(U=1\) and \(V\) is nonconstant.  Since
\(\deg C+\deg V\le\deg G\le2\), both \(C\) and \(V\) are linear.  The two
factors occur with multiplicity one in \(G\), while \(F\) has only the
support \(C\).  Hence, on writing \(L=C\), there are nonzero integers
\(\gamma,\delta\) and \(r\in\{1,2\}\) such that
\[
  F=\gamma L^r,\qquad G=\delta LV.
\]
When \(V\) is intersective,
\cref{thm:simultaneous-dominance} applied to the factor \(V\mid G\) proves
the property.  Otherwise choose a prime \(p\) dividing the leading
coefficient of \(V\).  Primitivity makes every \(V(n)\) a \(p\)-adic unit.
Use the lacunary test set
\[
  \mathcal A=\{p^{3^j}:j\ge1\}.
\]
For \(X\ge0\), write
\[
  N(X)=\#\{j\ge1:3^j\le X\}.
\]
At every nonzero value, with \(t=v_p(L(n))\), the two divisor counts are
\[
  N\bigl(v_p(\gamma)+rt\bigr)
  \quad\text{and}\quad
  N\bigl(v_p(\delta)+t\bigr).
\]
Their difference is uniformly bounded in \(t\): the two affine arguments,
after adding one, have a ratio bounded above and below, while
\(N(X)=\log_3(X+1)+O(1)\).  A zero of \(L\) is common to both polynomials,
and \(V\) has no integer zero.  Hence the extended difference is bounded and
the property fails.  The case in which only \(U\) is nonconstant is
symmetric.

In the remaining case \(C,U,V\) are all nonconstant.  The inequalities
\[
  \deg C+\deg U\le\deg F\le2,
  \qquad
  \deg C+\deg V\le\deg G\le2
\]
force all three pairwise coprime support factors to be linear.  Equality in
both inequalities also forces every factor multiplicity on the corresponding
side to be one, so this is the case
\(r_L=r_U=s_L=s_V=1\) of \cref{thm:linear-defects}.

These cases exhaust the possible supports in degree at most two.  The
criterion uses only \(C,U,V\), and intersectivity is unchanged by contents
and multiplicities by \cref{lem:radical-intersective}.  Applying the same
criterion to \(\radQ(F),\radQ(G)\) proves the final equivalence.
\end{proof}

\begin{proof}[Proof of \cref{cor:normalized-linear}]
If \(a=b\), the two polynomials are identical.  Otherwise apply
\cref{thm:linear-defects} with \(L=x\), \(U=ax+1\), \(V=bx+1\), and all
four exponents equal to one.
\end{proof}

When \(\gcd(|a|,|b|)=1\) and \(|a|,|b|>1\), the unique factors \(ax+1\)
and \(bx+1\) in \cref{cor:normalized-linear} are not intersective, whereas
their product is intersective: for each prime power, at least one of \(a,b\)
is invertible, and the Chinese remainder theorem combines the corresponding
linear roots.
Thus neither half of \cref{cor:sandwich} alone implies the conclusion.  The
choice \(a=2,b=3\) decides the basic model
\[
  F(x)=x(2x+1),\qquad G(x)=x(3x+1).
\]

\begin{question}\label[question]{q:exceptional-boundary}
Give a finite local--Galois criterion for \(\cP(F,G)\) when \(UV\) is
prime-covering but neither \(U\) nor \(V\) is intersective.  Beyond
\cref{cor:degree-two}, decide higher-degree split pairs, beginning with four
linear support factors or with nonlinear support factors.
\end{question}

Any complete answer must combine the permutation-group covering of the roots
of \(UV\) with the \(p\)-adic branches at the bad primes.  Prime density alone
cannot distinguish \cref{ex:236} from the intersective example following it;
multiplicities alone are also insufficient because prime test sets ignore
them.

\section{Acknowledgments and AI Disclosure}
  The author thanks Yuchen Ding for introducing this issue to him and for providing certain assistance in terms of writing and literature. The author thanks his graduate supervisor Huixi Li for providing him with many useful suggestions regarding the revision of the article.
  
  In a manuscript written around last November, the author developed the rudiments of \textit{intersective polynomial} and \textit{prime-covering}, and excluded some cases where \(\cP(F,G)\) did not hold. However, at that time, no AI tools were utilized. The author made multiple iterations on the original draft using ChatGPT 5.6Sol, completing the collection of literature, the writing of the article, including the establishment of the main lemma~\ref{lem:uniform-saturation}. Nevertheless, the final proofs in this paper have been reviewed by the human author, the author is responsible for this.

\begin{thebibliography}{10}

\bibitem{BerendBilu}
Daniel Berend and Yuri Bilu, \emph{Polynomials with roots modulo every
  integer}, Proceedings of the American Mathematical Society \textbf{124}
  (1996), no.~6, 1663--1671.

\bibitem{BergelsonLeibmanLesigne}
Vitaly Bergelson, Alexander Leibman, and Emmanuel Lesigne, \emph{Intersective
  polynomials and the polynomial {S}zemer{\'e}di theorem}, Advances in
  Mathematics \textbf{219} (2008), no.~1, 369--388.

\bibitem{DingLinear}
Yuchen Ding, \emph{Note on a problem of {S}{\'a}rk{\H{o}}zy on multiplicative
  representation functions}, European Journal of Combinatorics \textbf{132}
  (2026), Paper No.~104268, Part B.

\bibitem{DingProblem26}
Yuchen Ding, \emph{On problem 26 of {S}{\'a}rk{\H{o}}zy's list}, Periodica
  Mathematica Hungarica \textbf{92} (2026), no.~1, 148--150.

\bibitem{ElsholtzKlahnTechnau}
Christian Elsholtz, Benjamin Klahn, and Marc Technau, \emph{On polynomials
  with roots modulo almost all primes}, Acta Arithmetica \textbf{205} (2022),
  no.~3, 251--263.

\bibitem{HalberstamRichert}
Heini Halberstam and Hans-Egon Richert, \emph{Sieve methods}, London
  Mathematical Society Monographs, vol.~4, Academic Press, London--New York,
  1974.

\bibitem{HamelLyallRice}
Mariah Hamel, Neil Lyall, and Alex Rice, \emph{Improved bounds on
  {S}{\'a}rk{\H{o}}zy's theorem for quadratic polynomials}, International
  Mathematics Research Notices \textbf{2013} (2013), no.~8, 1761--1782.

\bibitem{KamaeMendesFrance}
Teturo Kamae and Michel Mend{\`e}s~France, \emph{Van der {C}orput's difference
  theorem}, Israel Journal of Mathematics \textbf{31} (1978), no.~3--4,
  335--342.

\bibitem{Le2010}
Th\'ai~Ho\`ang L{\^e}, \emph{Intersective polynomials and the primes}, Journal
  of Number Theory \textbf{130} (2010), no.~8, 1705--1717.

\bibitem{Lucier2006}
Jason Lucier, \emph{Intersective sets given by a polynomial}, Acta Arithmetica
  \textbf{123} (2006), no.~1, 57--95.

\bibitem{MishraMinimal}
Bhawesh Mishra, \emph{Minimally intersective polynomials with arbitrarily
  long factorization}, INTEGERS \textbf{22} (2022), Paper No.~A15.

\bibitem{MishraQuadratic}
Bhawesh Mishra, \emph{Polynomials consisting of quadratic factors with roots
  modulo any positive integer}, The American Mathematical Monthly
  \textbf{129} (2022), no.~2, 178--182.

\bibitem{Mishra2024}
Bhawesh Mishra, \emph{Polynomials with factors of the form \((x^q-a)\) with
  roots modulo every integer}, Communications in Algebra \textbf{52} (2024),
  no.~11, 4946--4954.

\bibitem{Sarkozy2001}
Andr{\'a}s S{\'a}rk{\H{o}}zy, \emph{Unsolved problems in number theory},
  Periodica Mathematica Hungarica \textbf{42} (2001), no.~1--2, 17--35.

\end{thebibliography}

\end{document}
